\section{Objetivos}

\subsection{Objetivo general}

Desarrollar una plataforma distribuida de difusión universitaria basada en una arquitectura \textit{publish-subscribe} (pub-sub) que permita a estudiantes, profesores y coordinación académica publicar y recibir notificaciones en tiempo real mediante canales de suscripción orientados a la alta usabilidad y al descubrimiento abierto de información dentro de la comunidad universitaria.

La plataforma será desarrollada como una aplicación móvil, debido a que los dispositivos móviles representan actualmente el principal medio de acceso a herramientas digitales y comunicación entre estudiantes.

\subsection{Objetivos específicos}

A partir del objetivo general, se definen los siguientes objetivos específicos, cada uno asociado a un conjunto de actividades técnicas y a las tecnologías seleccionadas en el marco teórico \cite{tanenbaum2023distributed, tarkoma2012publish}:

\begin{enumerate}
    \item \textbf{Diseñar una arquitectura distribuida} basada en el modelo \textit{publish-subscribe} para la gestión escalable de eventos y notificaciones en tiempo real, utilizando \textit{Spring Boot} para el desarrollo del \textit{backend}, \textit{Apache Kafka} para la distribución de eventos y \textit{Docker} para la contenerización y despliegue de servicios \cite{kafka2017streaming, spring2023boot}.
    
    \item \textbf{Implementar un sistema de canales y suscripciones} que permita la difusión y descubrimiento de información académica, cultural y comunitaria, utilizando \textit{PostgreSQL} para el almacenamiento de datos, \textit{Redis} para optimización y manejo de información temporal \cite{postgresql2024docs}, y \textit{JWT} para autenticación segura de usuarios \cite{jwt2015rfc}.
    
    \item \textbf{Desarrollar una aplicación móvil multiplataforma} utilizando \textit{Flutter} para la interfaz móvil y el envío de notificaciones push, priorizando principios de usabilidad, accesibilidad y rapidez de interacción orientados al entorno universitario \cite{flutter2024docs, amexcomp2020ihc}.
\end{enumerate}

\subsection{Alcance y entregables}

Al finalizar el proyecto, se obtendrán los siguientes entregables:

\begin{itemize}
    \item Un \textit{backend} funcional desplegable mediante contenedores Docker, que exponga una API REST para la gestión de usuarios, canales, publicaciones y suscripciones.
    \item Un sistema de eventos basado en Apache Kafka que gestione las notificaciones en tiempo real.
    \item Una base de datos PostgreSQL con el modelo relacional para almacenamiento persistente, complementada con Redis para caché.
    \item Una aplicación móvil desarrollada en Flutter para Android e iOS que permita:
    \begin{itemize}
        \item Registro e inicio de sesión mediante autenticación JWT.
        \item Creación y administración de canales temáticos.
        \item Suscripción a canales públicos y recepción de notificaciones push.
        \item Publicación de avisos dentro de los canales.
        \item Búsqueda y descubrimiento de canales por categoría o palabra clave.
    \end{itemize}
    \item Documentación técnica del sistema (diagramas de arquitectura, manual de instalación, guía de usuario).
    \item Reporte final del proyecto que incluya resultados de pruebas de usabilidad y rendimiento.
\end{itemize}

\subsection{Relación con el marco teórico}

Cada objetivo específico se fundamenta en los conceptos expuestos en el marco teórico. La siguiente tabla muestra dicha correspondencia:

\begin{table}[H]
\centering
\caption{Correspondencia entre objetivos específicos y fundamentos teóricos}
\label{tab:objetivos_marco}
\small
\begin{tabular}{|p{4cm}|p{4cm}|}
\hline
\textbf{Objetivo específico} & \textbf{Fundamento teórico} \\
\hline
1. Arquitectura distribuida pub-sub & Modelo publish-subscribe \cite{tarkoma2012publish}, sistemas distribuidos \cite{tanenbaum2023distributed}, Spring Boot \cite{spring2023boot}, Kafka \cite{kafka2017streaming} \\
\hline
2. Canales, suscripciones y autenticación & Persistencia con PostgreSQL \cite{postgresql2024docs}, caché con Redis, autenticación JWT \cite{jwt2015rfc} \\
\hline
3. Aplicación móvil y usabilidad & Flutter \cite{flutter2024docs}, principios de Interacción Humano-Computadora \cite{amexcomp2020ihc} \\
\hline
\end{tabular}
\end{table}

\subsection{Indicadores de éxito}

Para evaluar el cumplimiento de los objetivos, se definen los siguientes indicadores:

\begin{itemize}
    \item \textbf{Funcionalidad completa:} La plataforma debe implementar todas las características listadas en el alcance (canales, suscripciones, notificaciones push, autenticación).
    \item \textbf{Rendimiento:} El sistema debe ser capaz de manejar altos volúmenes de tráfico con baja latencia.
    \item \textbf{Usabilidad:} La aplicación móvil debe obtener una puntuación superior a 70 en la escala SUS (System Usability Scale) tras pruebas con al menos 15 estudiantes.
    \item \textbf{Despliegue reproducible:} El sistema completo debe poder levantarse en un entorno limpio usando únicamente el comando \texttt{docker-compose up}.
\end{itemize}

Estos objetivos e indicadores guiarán el desarrollo incremental descrito en la metodología (sección 6).